\section{Результаты моделирования}

\subsection{Параметры модели}

Численное моделирование выполнено для радиального подшипника скольжения геометрии~\varGeomLetter{} ($R = \varRmm$~мм, $c = \varCmkm$~мкм, $L = \varLmm$~мм) с \varDimpleTypeGenitive{} микрорельефом. Диаметр вала составляет $D = 2R = 100$~мм. Исходные данные для расчёта представлены в таблице~\ref{tab:input_data}.

\begin{table}[H]
\small
\caption{Исходные данные для расчёта}
\label{tab:input_data}
\begin{tabularx}{\textwidth}{|X|c|c|}
\hline
\centering\textbf{Параметр} & \textbf{Значение} & \textbf{Размерность} \\
\hline
Радиус вала $R$ & \varRmm & мм \\
\hline
Радиальный зазор $c$ & \varCmkm & мкм \\
\hline
Длина подшипника $L$ & \varLmm & мм \\
\hline
Относительная длина $L/D$ & \varLD & -- \\
\hline
Полуось по оси $Z$, $a$ & \varDimpleAmm & мм \\
\hline
Полуось по оси $\varphi$, $b$ & \varDimpleBmm & мм \\
\hline
Глубина углубления $h_p$ & \varDimpleHpmkm & мкм \\
\hline
Относительная глубина $h_p/c$ & \varDimpleHpRatio & -- \\
\hline
Число углублений по окружности $N_\varphi$ & \varNphi & шт. \\
\hline
Число рядов по длине $N_Z$ & \varNZ & шт. \\
\hline
Область текстурирования & \varPhiRange & -- \\
\hline
Частота вращения $n$ & \varN & об/мин \\
\hline
Вязкость смазки $\eta$ & \varEta & Па$\cdot$с \\
\hline
Диапазон эксцентриситета $\varepsilon$ & 0{,}05--0{,}80 & -- \\
\hline
Размер сетки & $500 \times 500$ & -- \\
\hline
Параметр релаксации $\omega$ & 1{,}5 & -- \\
\hline
Критерий сходимости $\delta$ & $< 10^{-5}$ & -- \\
\hline
Макс. число итераций & 50000 & -- \\
\hline
\end{tabularx}
\end{table}

Линейная скорость поверхности вала при заданной частоте вращения равна $U = \pi D n / 60 = \pi \cdot 0{,}100 \cdot 2980 / 60 \approx 15{,}60$~м/с. Расчёт проведён для 15 значений эксцентриситета в указанном диапазоне.

\subsection{Распределение зазора и давления}

На рисунке~\ref{fig:pressure_2d} представлено распределение безразмерного давления $P(\varphi)$ в среднем сечении ($Z = 0$) при $\varepsilon = 0{,}6$ для гладкого и текстурированного подшипников.

\begin{figure}[H]
\centering
\includegraphics[width=\textwidth]{figures/pressure_2d.png}
\caption{Распределение давления $P(\varphi)$ при $Z=0$ для гладкого и текстурированного подшипников ($\varepsilon = 0{,}6$)}
\label{fig:pressure_2d}
\end{figure}

Для гладкого подшипника давление имеет единственный максимум в районе минимального зазора ($\varphi \approx \pi$). По мере сужения зазора (от $\varphi = 0$ до $\varphi \approx \pi$) давление монотонно нарастает, а в расширяющейся части обращается в нуль вследствие кавитации. Положительное давление занимает примерно половину окружности подшипника. Наличие цилиндрического микрорельефа вносит характерные возмущения в кривую давления. Главной отличительной чертой цилиндрического профиля является постоянная глубина внутри границ углубления и резкие перепады зазора на его кромках. Это порождает интенсивные локальные градиенты давления: выраженный пик на входной кромке, где поток встречает ступенчатое расширение, и спад на выходной кромке. Благодаря кавитационному ограничению ($P \geq 0$) суммарный вклад каждого углубления остаётся положительным, а общий уровень давления текстурированного подшипника оказывается выше, чем у гладкого.

На рисунке~\ref{fig:clearance_2d} показан профиль зазора $H(\varphi)$ в среднем сечении при $\varepsilon = 0{,}6$.

\begin{figure}[H]
\centering
\includegraphics[width=\textwidth]{figures/clearance_2d.png}
\caption{Распределение зазора $H(\varphi)$ при $Z=0$ для подшипника с \varDimpleTypeGenitive{} углублениями ($\varepsilon = 0{,}6$)}
\label{fig:clearance_2d}
\end{figure}

Зазор изменяется по закону $H = 1 + \varepsilon\cos\varphi$, а в зонах текстурирования к нему добавляются ступенчатые углубления постоянной глубины. В безразмерных единицах глубина каждого углубления составляет $\Delta H_{\max} = h_p/c = 0{,}1$. Переходы на кромках углублений носят скачкообразный характер, что типично для цилиндрического профиля и отличает его от плавных эллипсоидальных и параболических форм.

Трёхмерные распределения зазора и давления при $\varepsilon = 0{,}6$ приведены на рисунках~\ref{fig:clearance_3d} и~\ref{fig:pressure_3d}.

\begin{figure}[H]
\centering
\includegraphics[width=\textwidth]{figures/clearance_3d.png}
\caption{Трёхмерное распределение зазора $H(\varphi, Z)$ ($\varepsilon = 0{,}6$)}
\label{fig:clearance_3d}
\end{figure}

На рисунке~\ref{fig:clearance_3d} отчётливо виден регулярный массив цилиндрических углублений: 8 по окружности и 11 рядов по длине подшипника, расположенных в зоне $90^\circ - 270^\circ$. Каждое углубление представляет собой область постоянной дополнительной глубины с чётко выраженными вертикальными стенками.

\begin{figure}[H]
\centering
\includegraphics[width=\textwidth]{figures/pressure_3d.png}
\caption{Трёхмерное распределение давления $P(\varphi, Z)$ ($\varepsilon = 0{,}6$)}
\label{fig:pressure_3d}
\end{figure}

На рисунке~\ref{fig:pressure_3d} наряду с основным максимумом давления, обусловленным клиновой геометрией зазора, наблюдаются интенсивные локальные возмущения на каждом углублении. Резкие кромки цилиндрического профиля формируют сильные градиенты давления, что отличает данный тип от плавных профилей, где возмущения более сглажены. На входной кромке углубления зазор скачкообразно увеличивается, создавая выраженный подъём давления~--- микрогидродинамический эффект. На выходной кромке происходит спад давления, однако асимметрия вклада за счёт кавитации обеспечивает положительный суммарный эффект.

Анализ трёхмерных распределений позволяет выделить следующие закономерности. Максимальное давление формируется в среднем сечении подшипника ($Z = 0$), а к торцам ($Z = \pm 1$) давление убывает до нуля вследствие торцевых утечек. Эффект текстурирования наиболее выражен вблизи среднего сечения, где фоновое давление достигает наибольших значений. Для подшипника геометрии~B, отличающегося увеличенными размерами ($D = 100$~мм, $L = 80$~мм), абсолютные значения нагрузочной способности существенно выше по сравнению с геометрией~A, что обусловлено большей площадью несущей поверхности.

\subsection{Сравнительный анализ интегральных показателей}

Интегральные характеристики гладкого подшипника и подшипника с цилиндрическими углублениями для выбранных значений эксцентриситета приведены в таблицах~\ref{tab:smooth_results} и~\ref{tab:textured_results}.

\begin{table}[H]
\small
\centering
\caption{Характеристики гладкого подшипника}
\label{tab:smooth_results}
\begin{tabular}{|c|c|c|c|}
\hline
$\boldsymbol{\varepsilon}$ & $\boldsymbol{F}$\textbf{, Н} & $\boldsymbol{\mu}$ & $\boldsymbol{Q}$\textbf{, л/с} \\
\hline
0{,}05 & 551{,}5 & 0{,}2245 & 0{,}6852 \\
\hline
0{,}16 & 1775{,}8 & 0{,}0707 & 0{,}6843 \\
\hline
0{,}26 & 3141{,}8 & 0{,}0412 & 0{,}6825 \\
\hline
0{,}37 & 4805{,}2 & 0{,}0282 & 0{,}6801 \\
\hline
0{,}48 & 6984{,}6 & 0{,}0207 & 0{,}6772 \\
\hline
0{,}59 & 10124{,}4 & 0{,}0157 & 0{,}6744 \\
\hline
0{,}69 & 15347{,}9 & 0{,}0118 & 0{,}6717 \\
\hline
0{,}80 & 25964{,}3 & 0{,}0085 & 0{,}6698 \\
\hline
\end{tabular}
\end{table}

\begin{table}[H]
\small
\centering
\caption{Характеристики подшипника с \varDimpleTypeGenitive{} углублениями}
\label{tab:textured_results}
\begin{tabular}{|c|c|c|c|}
\hline
$\boldsymbol{\varepsilon}$ & $\boldsymbol{F}$\textbf{, Н} & $\boldsymbol{\mu}$ & $\boldsymbol{Q}$\textbf{, л/с} \\
\hline
0{,}05 & 753{,}7 & 0{,}1633 & 0{,}6904 \\
\hline
0{,}16 & 2002{,}2 & 0{,}0623 & 0{,}6894 \\
\hline
0{,}26 & 3449{,}1 & 0{,}0372 & 0{,}6876 \\
\hline
0{,}37 & 5240{,}8 & 0{,}0256 & 0{,}6851 \\
\hline
0{,}48 & 7671{,}2 & 0{,}0187 & 0{,}6822 \\
\hline
0{,}59 & 11435{,}1 & 0{,}0137 & 0{,}6792 \\
\hline
0{,}69 & 18355{,}3 & 0{,}0097 & 0{,}6763 \\
\hline
0{,}80 & 36204{,}1 & 0{,}0060 & 0{,}6739 \\
\hline
\end{tabular}
\end{table}

Нагрузочная способность $F$ нелинейно возрастает с увеличением эксцентриситета для обоих вариантов подшипника. Подшипник с цилиндрическими углублениями обладает более высокой нагрузочной способностью при всех значениях $\varepsilon$. При максимальном эксцентриситете ($\varepsilon = 0{,}80$) эффект текстурирования проявляется наиболее интенсивно: нагрузочная способность текстурированного подшипника (36204{,}1~Н) превышает гладкий аналог (25964{,}3~Н) на 39{,}4\%. Высокая эффективность цилиндрического профиля обусловлена тем, что резкие кромки углублений создают сильные локальные перепады давления, формирующие выраженный микрогидродинамический подъём.

Коэффициент трения $\mu$ монотонно убывает при увеличении $\varepsilon$ и для текстурированного подшипника принимает меньшие значения во всём исследованном диапазоне. Расход смазки $Q$ слабо зависит от эксцентриситета и у текстурированного подшипника незначительно выше вследствие увеличения эффективного объёма смазочного слоя в зонах углублений.

Для количественной оценки эффективности текстурирования в таблице~\ref{tab:comparison} выполнено сопоставление характеристик при $\varepsilon = 0{,}6$.

\begin{table}[H]
\small
\centering
\caption{Сравнение характеристик при $\varepsilon = 0{,}6$}
\label{tab:comparison}
\begin{tabular}{|l|c|c|c|}
\hline
\textbf{Характеристика} & \textbf{Гладкий} & \textbf{С углублениями} & $\boldsymbol{\Delta}$\textbf{, \%} \\
\hline
Нагрузка $F$, Н & 10685{,}6 & 12080{,}2 & $+13{,}05$ \\
\hline
Коэфф. трения $\mu$ & 0{,}01507 & 0{,}01319 & $-12{,}47$ \\
\hline
Расход $Q$, л/с & 0{,}6739 & 0{,}6788 & $+0{,}72$ \\
\hline
\end{tabular}
\end{table}

Нанесение цилиндрических углублений обеспечивает прирост несущей способности подшипника на 13{,}05\% за счёт микрогидродинамического эффекта на резких кромках. Коэффициент трения снижается на 12{,}47\% благодаря увеличению средней толщины смазочного слоя в зоне углублений. Расход смазки возрастает лишь на 0{,}72\%, что свидетельствует о минимальном влиянии текстурирования на торцевые утечки. Следует отметить, что для подшипника геометрии~B с увеличенными размерами ($D = 100$~мм) абсолютные значения нагрузочной способности существенно выше, чем для геометрии~A, однако относительный эффект текстурирования при тех же параметрах углублений выражен умереннее.

На рисунках~\ref{fig:load_vs_eps}--\ref{fig:flow_vs_eps} представлены зависимости интегральных характеристик от эксцентриситета для обоих вариантов подшипника.

\begin{figure}[H]
\centering
\includegraphics[width=\textwidth]{figures/load_vs_epsilon.png}
\caption{Нагрузочная способность $F(\varepsilon)$ для гладкого и текстурированного подшипников}
\label{fig:load_vs_eps}
\end{figure}

Кривая $F(\varepsilon)$ текстурированного подшипника располагается выше кривой гладкого подшипника во всём диапазоне эксцентриситетов. С увеличением $\varepsilon$ абсолютная разница между кривыми нарастает: при $\varepsilon = 0{,}05$ она составляет около 200~Н, тогда как при $\varepsilon = 0{,}80$ достигает более 10000~Н. Нелинейный характер зависимости обусловлен степенной связью давления с обратной толщиной зазора ($P \sim h^{-3}$), вследствие чего при малых зазорах вклад каждого углубления многократно усиливается.

\begin{figure}[H]
\centering
\includegraphics[width=\textwidth]{figures/friction_vs_epsilon.png}
\caption{Коэффициент трения $\mu(\varepsilon)$ для гладкого и текстурированного подшипников}
\label{fig:friction_vs_eps}
\end{figure}

Коэффициент трения текстурированного подшипника ниже гладкого при всех значениях $\varepsilon$. Обе кривые монотонно убывают, асимптотически приближаясь к нулю при $\varepsilon \to 1$. Наибольший абсолютный выигрыш наблюдается при малых эксцентриситетах: при $\varepsilon = 0{,}05$ разница составляет $\Delta\mu = 0{,}061$, тогда как при $\varepsilon = 0{,}80$~--- лишь $\Delta\mu = 0{,}003$. В относительном выражении снижение трения сохраняется на уровне 12--30\% во всём диапазоне эксцентриситетов.

\begin{figure}[H]
\centering
\includegraphics[width=\textwidth]{figures/flow_vs_epsilon.png}
\caption{Расход смазки $Q(\varepsilon)$ для гладкого и текстурированного подшипников}
\label{fig:flow_vs_eps}
\end{figure}

Расход смазки текстурированного подшипника незначительно превышает расход гладкого. Разница составляет менее 1\% и практически не зависит от эксцентриситета, что указывает на слабое влияние цилиндрических углублений на торцевые утечки.

Таким образом, текстурирование цилиндрическими углублениями обеспечивает одновременное повышение нагрузочной способности и снижение коэффициента трения при минимальном росте расхода смазки. Эффект усиливается при увеличении эксцентриситета, что представляет практическую ценность для тяжелонагруженных подшипниковых узлов.

\subsection*{Выводы по главе 3}
\phantomsection
\addcontentsline{toc}{subsection}{Выводы по главе 3}

Выполнено численное моделирование характеристик радиального подшипника скольжения геометрии~B с цилиндрическим микрорельефом для 15 значений эксцентриситета. Получены поля давления и зазора, определены зависимости нагрузочной способности, коэффициента трения и расхода смазки от эксцентриситета. Сравнение при $\varepsilon = 0{,}6$ выявило прирост нагрузочной способности на $13{,}05$\%, снижение коэффициента трения на $12{,}47$\% и рост расхода на $0{,}72$\%, что подтверждает эффективность цилиндрического текстурирования для подшипника увеличенных размеров.
